\section{Cross-Network Comparative Analysis}


Zoo's tokenomics must be contextualized within the three-tier ecosystem:

\subsection{Lux (L0) - Genesis NFT Controlled Release}

\textbf{Distribution Philosophy:} Gradual decentralization with quality control

\textbf{Mechanism:}
\begin{itemize}
\item Genesis NFT sale (5,000 NFTs at 0.1-10 ETH depending on rarity)
\item Each NFT unlocks 200K-10M $\Lux$ tokens (variable by tier)
\item Vesting: 20\% immediate, 80\% over 4 years
\item Staking requirement: NFT holders must stake 50\%+ of tokens to maintain benefits
\end{itemize}

\textbf{Rationale:} Base consensus layer requires stability. NFT sale ensures committed early adopters with financial stake. Gradual release prevents governance attacks during bootstrapping.

\textbf{Outcomes (as of Oct 2025):}
\begin{itemize}
\item 4,287 NFTs sold (86\% sell-through)
\item 220B $\Lux$ unlocked (22\% of community allocation)
\item Avg holder stake: 68\% (exceeds 50\% requirement)
\item Network uptime: 99.97\%
\end{itemize}

\subsection{Hanzo (L1) - Fair Launch Proof-of-Compute}

\textbf{Distribution Philosophy:} Meritocratic compute contribution

\textbf{Mechanism:}
\begin{itemize}
\item No pre-mine, no ICO, no VC allocation
\item Proof-of-compute mining: Submit AI/ML training jobs to earn $\Hanzo$
\item Block reward: 50,000 $\Hanzo$ per block (halves every 1,051,200 blocks ≈ 4 years)
\item Difficulty adjustment: Every 2,016 blocks based on network hashrate
\end{itemize}

\textbf{Rationale:} Compute layer serves as infrastructure for AI workloads. Mining ensures participants contribute actual GPU/TPU resources before receiving tokens.

\textbf{Outcomes (as of Oct 2025):}
\begin{itemize}
\item 84B $\Hanzo$ mined (8.4\% of minable supply)
\item 1,247 active miners
\item Network compute: 2.4 exaFLOPS (2.4 × 10\textsuperscript{18} FLOPS)
\item Gini coefficient: 0.73 (moderate concentration - early miners advantage)
\end{itemize}

\subsection{Zoo (L2) - 100\% Community Airdrop}

\textbf{Distribution Philosophy:} Instant democratic ownership

\textbf{Mechanism:}
\begin{itemize}
\item One-time airdrop based on contribution scoring (data, code, engagement)
\item No sales, no mining, no vesting
\item Claim period: 90 days (unclaimed tokens burned)
\end{itemize}

\textbf{Rationale:} Training-free AI requires community-generated semantic experiences. Airdrop incentivizes immediate high-quality contribution. Specialized layer (not base consensus) can tolerate instant decentralization.

\textbf{Outcomes (as of Oct 2025):}
\begin{itemize}
\item 720B $\KEEPER$ claimed (72\% of airdrop)
\item 62,885 claimants (72\% claim rate)
\item Gini coefficient: 0.61 (most equitable of three networks)
\item Experience contributions: 12.4M examples (6 weeks post-airdrop)
\end{itemize}

\subsection{Comparative Table}

\begin{table}[h]
\centering
\small
\begin{tabular}{lcccc}
\toprule
\textbf{Metric} & \textbf{Lux (L0)} & \textbf{Hanzo (L1)} & \textbf{Zoo (L2)} & \textbf{Industry Avg} \\
\midrule
Total Supply & 2T & 2T & 2T & 10B-100B \\
DAO Allocation & 1T (50\%) & 1T (50\%) & 1T (50\%) & 20-30\% \\
Distribution & NFT sale & Mining & Airdrop & ICO/VC \\
Capital Raised & \$15M & \$0 & \$0 & \$50-500M \\
Gini Coefficient & 0.68 & 0.73 & 0.61 & 0.80-0.90 \\
Unique Holders & 4,287 & 1,247 & 62,885 & 10K-100K \\
Fair Launch? & Partial & Yes & Yes & No \\
Inflation & 0\% & 2-4\% & 0\% & 2-10\% \\
Governance & NFT-weighted & Mining-weighted & Quadratic & Token-weighted \\
\bottomrule
\end{tabular}
\caption{Comparative tokenomics across Lux-Hanzo-Zoo ecosystem}
\label{tab:cross_network}
\end{table}

\subsection{Validator Requirements Across Networks}

Each network implements distinct validator staking requirements optimized for its role:

\begin{table}[h]
\centering
\begin{tabular}{lccc}
\toprule
\textbf{Requirement} & \textbf{Lux (L0)} & \textbf{Hanzo (L1)} & \textbf{Zoo (L2)} \\
\midrule
Min Stake & 1M $\Lux$ & 1 AI & 1,000 $\KEEPER$ \\
Validator Count & 100 (fixed) & Unlimited & Unlimited \\
Validator Allocation & 1B/each & None & None \\
Unlock Schedule & 100 years & N/A & N/A \\
Entry Barrier & High (security) & 1 AI (self-mined) & Low (participation) \\
Validation Method & PoS + Genesis & PoW (compute) & PoAI (Proof of AI) \\
\bottomrule
\end{tabular}
\caption{Validator staking requirements across Lux-Hanzo-Zoo}
\label{tab:validator_comparison}
\end{table}

\textbf{Lux (L0) - High Security Model:}
\begin{itemize}
\item 100 genesis validators each allocated 1 billion $\Lux$ tokens
\item Tokens unlock linearly over 100 years (10M/year per validator)
\item 1M $\Lux$ minimum stake ensures serious commitment (\$20M+ at \$20/token)
\item Fixed validator set during bootstrapping phase (first 2 years)
\item Rationale: Base consensus layer requires maximum security and stability
\end{itemize}

\textbf{Hanzo (L1) - Self-Mined Compute Participation:}
\begin{itemize}
\item 1 AI token minimum stake (self-mined on any device - laptop, phone, etc.)
\item Mine your first 1 AI token locally, then become a Hanzo validator
\item Participate in HMM (Hamiltonian Market Maker) compute market consensus
\item Provide compute resources on Blackwell GPU infrastructure
\item Fully private public AI chain with permissionless participation
\item Validator rewards scale with compute contribution (FLOPS delivered)
\item Rationale: Low barrier (just mine 1 token) ensures participation commitment while maintaining accessibility
\end{itemize}

\textbf{Zoo (L2) - Proof of AI (PoAI) Semantic Validation:}
\begin{itemize}
\item 1,000 $\KEEPER$ minimum stake (approx \$20 at \$0.02/token)
\item \textbf{PoAI weighting}: Validation power scaled by LLM experience contributions
\item Experience Ledger participation: Validators curate and vote on semantic experiences
\item Semantic optimization rewards: Earn additional rewards for high-quality experience submissions
\item Byzantine-robust aggregation: DAO voting (66\% threshold) prevents malicious experiences
\item Enables broad validator participation across 142 countries
\item Delegation available for non-validators
\item Rationale: AI specialization layer benefits from diverse semantic contributions and cultural perspectives
\end{itemize}

This tiered approach balances security (Lux), accessibility (Hanzo), and democratic participation (Zoo) across the ecosystem.

\subsection{Interoperability and Value Flow}

The three networks exhibit symbiotic value flows:

\begin{itemize}
\item \textbf{Lux → Hanzo}: Base chain security inherits to compute layer
\item \textbf{Hanzo → Zoo}: GPU compute resources power Zoo inference/training
\item \textbf{Zoo → Hanzo}: Experience library accessible to general Hanzo workloads
\item \textbf{Zoo → Lux}: DAO governance proposals can request Lux treasury grants
\end{itemize}

\textbf{Token Bridge Economics:}

Users can bridge tokens across networks:
\begin{itemize}
\item $\Lux$ $\leftrightarrow$ $\Hanzo$: 1:1 peg, 0.1\% bridge fee
\item $\Hanzo$ $\leftrightarrow$ $\KEEPER$: 1:1 peg, 0.1\% bridge fee
\item $\Lux$ $\leftrightarrow$ $\KEEPER$: 1:1 peg, 0.2\% bridge fee (two-hop)
\end{itemize}

Bridge fees split:
\begin{itemize}
\item 60\% → Bridge validators (stake $\Lux$ + $\Hanzo$ + $\KEEPER$)
\item 40\% → DAO treasuries (proportional to hop distance)
\end{itemize}

This creates \textbf{economic interdependence} - networks benefit from each other's growth.

